% !PRINTED 556
\weberpage{556}

If one proceeds in this way with all primes and sets
\[
S'_{10} + S'_{01} - 2\,S'_{11} = T',
\]
then it follows that
\[
T' = T\,\Pi\left(1 - \frac{1}{p^2}\right),
\]
in which the product $\Pi$ extends over \gesperrt{all odd primes} $p$. Now, by development in
ascending powers of $p^{-2}$, one obtains:
\[
\begin{aligned}
\frac{1}{\Pi(1-p^{-2})} &= 1 + \frac{1}{3^2} + \frac{1}{5^2} + \frac{1}{7^2} + \frac{1}{9^2} + \cdots,\\
&= \frac{\pi^2}{8},
\end{aligned}
\]
and consequently
\[
T = \frac{\pi^2}{8}\,T',
\]
and by proceeding in the same way with the two other sums, there follows from (9):
\weq{10}{\begin{aligned}
\log f(\omega_1)\,f(\omega_2) &= \tfrac{1}{3}\log 2 + \frac{\beta\pi}{24}\left(S'_{10} + S'_{01} - 2\,S'_{11}\right), \\
\log f_1(\omega_1)\,f_1(\omega_2) &= \tfrac{1}{3}\log 2 + \frac{\beta\pi}{24}\left(S'_{10} - 2\,S'_{01} + S'_{11}\right), \\
\log f_2(\omega_1)\,f_2(\omega_2) &= \tfrac{1}{3}\log 2 + \frac{\beta\pi}{24}\left(-2\,S'_{10} + S'_{01} + S'_{11}\right).
\end{aligned}}

These formulae simplify essentially if one sets $\alpha = 0$, and hence
\[
\omega_1 = \omega_2 = \omega = i\,\beta
\]
that is, if one assumes a \gesperrt{purely imaginary} $\omega$. Then
\[
S'_{10} = \sum \frac{1}{x^2 - \omega^2 y^2} \qquad \begin{array}{l}x \text{ odd}, \ y \text{ even},\\xy \text{ relatively prime}.\end{array}
\]

If one separates off the term $x = \pm 1$, $y = 0$ and takes the remaining ones together four
at a time, then there results
\[
S'_{10} = 2 + 4\sum \frac{1}{x^2 - \omega^2 y^2},
\]
in which $x$, $y$ run through \gesperrt{positive values} only.

One proceeds in the same way with the other sums. If one then sets
%CONT>>
